%==============================================================================
% Chapitre 41 - L'étui
%==============================================================================

\section{Introduction}

L'étui est l'élément central de la cartouche moderne.

Souvent considéré comme un simple contenant, il joue pourtant un rôle
fondamental dans :

\begin{itemize}
    \item le fonctionnement de l'arme ;
    \item la sécurité ;
    \item la fiabilité ;
    \item la précision.
\end{itemize}

Il assure le maintien des différents composants de la cartouche et participe
activement au déroulement du tir.

\begin{aretenir}

L'étui n'est pas un simple récipient : il constitue un composant mécanique
essentiel de la cartouche.

\end{aretenir}

\begin{figure}[htbp]
\centering
\begin{tikzpicture}[scale=1.2]

% Corps principal
\draw[thick] (0,0) -- (0,5);
\draw[thick] (2,0) -- (2,5);

% Épaulement
\draw[thick] (0,5) -- (0.5,6);
\draw[thick] (2,5) -- (1.5,6);

% Collet
\draw[thick] (0.5,6) -- (0.5,7);
\draw[thick] (1.5,6) -- (1.5,7);

% Projectile
\draw[thick] (0.5,7) -- (1,8);
\draw[thick] (1.5,7) -- (1,8);

% Culot
\draw[thick] (-0.2,0) -- (2.2,0);

% Amorce
\draw[fill=white] (1,0.25) circle (0.15);

% Annotations
\draw[->] (3,7.5) -- (1,7.6);
\node[right] at (3,7.5) {Projectile};

\draw[->] (3,6.3) -- (1.5,6.5);
\node[right] at (3,6.3) {Collet};

\draw[->] (3,5.3) -- (1.7,5.5);
\node[right] at (3,5.3) {Épaulement};

\draw[->] (3,3.5) -- (2,3.5);
\node[right] at (3,3.5) {Corps};

\draw[->] (3,0.7) -- (2,0.3);
\node[right] at (3,0.7) {Culot};

\draw[->] (-2,0.5) -- (1,0.25);
\node[left] at (-2,0.5) {Amorce};

\end{tikzpicture}
\caption{Principales parties d'un étui à gorge de type bottleneck.}
\label{fig:parties_etui}
\end{figure}

\section{Fonctions principales}

L'étui remplit plusieurs fonctions simultanément :

\begin{itemize}
    \item contenir la charge propulsive ;
    \item maintenir le projectile ;
    \item recevoir l'amorce ;
    \item permettre le chargement ;
    \item assurer l'étanchéité lors du tir ;
    \item faciliter l'extraction après le tir.
\end{itemize}

Chaque fonction influence directement le fonctionnement de l'arme.

\section{Évolution historique}

Les premières armes à feu ne possédaient pas d'étui métallique.

La poudre et le projectile étaient chargés séparément.

L'apparition des cartouches métalliques au XIX\ieme\ siècle a permis :

\begin{itemize}
    \item d'accélérer le chargement ;
    \item d'améliorer la fiabilité ;
    \item d'augmenter les cadences de tir ;
    \item de supporter des pressions plus élevées.
\end{itemize}

Cette innovation a profondément transformé les armes à feu.

\section{Constitution générale}

Un étui moderne comporte généralement plusieurs zones distinctes :

\begin{itemize}
    \item le collet ;
    \item l'épaulement ;
    \item le corps ;
    \item le culot ;
    \item la gorge d'extraction.
\end{itemize}

Chacune possède une fonction spécifique.

\section{Le collet}

Le collet est la partie supérieure de l'étui.

Il assure le maintien du projectile.

Sa géométrie influence notamment :

\begin{itemize}
    \item la tension de collet ;
    \item la régularité du départ du projectile ;
    \item la précision.
\end{itemize}

\section{L'épaulement}

L'épaulement relie le corps de l'étui au collet.

Dans de nombreuses cartouches modernes, il sert également de référence pour
le positionnement dans la chambre.

On parle alors de prise de feuillure sur l'épaulement.

\section{Le corps}

Le corps constitue la plus grande partie du volume de l'étui.

Il contient principalement :

\begin{itemize}
    \item la poudre ;
    \item les gaz produits lors du tir.
\end{itemize}

Sa capacité interne influence directement les performances balistiques.

\section{Le culot}

Le culot est la partie arrière de l'étui.

Il supporte les pressions les plus importantes lors du tir.

Cette zone doit présenter une résistance mécanique élevée.

Elle comporte généralement :

\begin{itemize}
    \item le logement d'amorce ;
    \item la gorge d'extraction ;
    \item éventuellement un bourrelet.
\end{itemize}

\section{La gorge d'extraction}

La gorge d'extraction permet à l'extracteur de saisir l'étui après le tir.

Elle joue un rôle essentiel dans le fonctionnement des armes à répétition et
semi-automatiques.

\section{Le bourrelet}

Certaines cartouches possèdent un bourrelet dépassant du diamètre du corps.

On parle alors de cartouches à bourrelet.

Exemples :

\begin{itemize}
    \item .22 Long Rifle ;
    \item .38 Special ;
    \item .357 Magnum.
\end{itemize}

D'autres cartouches sont dites sans bourrelet apparent (\emph{rimless}).

\section{Matériaux}

La majorité des étuis modernes sont fabriqués en laiton.

Ce matériau présente plusieurs avantages :

\begin{itemize}
    \item bonne résistance mécanique ;
    \item excellente ductilité ;
    \item bonne résistance à la corrosion ;
    \item facilité de fabrication.
\end{itemize}

D'autres matériaux existent :

\begin{itemize}
    \item acier ;
    \item aluminium ;
    \item polymères.
\end{itemize}

\section{L'obturation}

Lors du tir, la pression augmente rapidement.

L'étui se dilate alors contre les parois de la chambre.

Ce phénomène est appelé obturation.

Il permet :

\begin{itemize}
    \item d'empêcher la fuite des gaz vers l'arrière ;
    \item de protéger le tireur ;
    \item d'améliorer le rendement énergétique.
\end{itemize}

\begin{aretenir}

L'obturation est l'une des fonctions les plus importantes de l'étui.

\end{aretenir}

\section{Retour élastique}

Après le tir, la pression diminue.

L'étui se contracte légèrement grâce à son élasticité.

Cette contraction facilite son extraction.

C'est l'une des raisons du succès du laiton.

\section{Capacité interne}

La capacité interne correspond au volume disponible pour la poudre.

Cette grandeur influence :

\begin{itemize}
    \item la pression ;
    \item la vitesse initiale ;
    \item les performances globales.
\end{itemize}

Deux étuis extérieurement identiques peuvent présenter des capacités
internes différentes.

\section{Tolérances de fabrication}

Les dimensions d'un étui sont définies avec précision.

Les tolérances influencent notamment :

\begin{itemize}
    \item la sécurité ;
    \item la précision ;
    \item la compatibilité avec les chambres.
\end{itemize}

\section{Étuis et rechargement}

Les étuis métalliques peuvent souvent être réutilisés.

Le rechargement consiste notamment à :

\begin{itemize}
    \item désamorcer ;
    \item recalibrer ;
    \item réamorcer ;
    \item recharger en poudre ;
    \item replacer un projectile.
\end{itemize}

Cette pratique sera étudiée dans un chapitre dédié.

\section{Influence sur la précision}

La qualité de l'étui influence :

\begin{itemize}
    \item la régularité des pressions ;
    \item la vitesse initiale ;
    \item la dispersion des tirs.
\end{itemize}

Les tireurs de précision accordent donc une attention particulière à ce
composant.

\section{Vers les amorces}

L'étui constitue l'ossature de la cartouche.

Cependant, sans système d'allumage, aucun tir ne pourrait avoir lieu.

Nous allons maintenant étudier l'amorce, premier élément actif de la chaîne
de mise à feu.

\section{À retenir}

\begin{aretenir}

\begin{itemize}
    \item L'étui regroupe les composants de la cartouche.
    \item Il assure le chargement, l'obturation et l'extraction.
    \item Le laiton est le matériau le plus utilisé.
    \item La capacité interne influence les performances balistiques.
    \item La qualité de l'étui participe à la précision.
\end{itemize}

\end{aretenir}
